% theme: sexual_reproduction_and_heredity
\MajorQuestion{有性生殖／遺伝の規則性と遺伝子}
\SetRowSpaceQanda

\Instruction{動物の有性生殖に関する次の問いに答えなさい。}
\QQ{雌と雄の生殖細胞の受精によって新しい個体をつくることを\NamedBlank{sexual_reproduction}という。}{有性生殖}
\QQ{雌の生殖細胞と、それをつくる器官をそれぞれ答えなさい。}{卵、卵巣}
\QQ{雄の生殖細胞と、それをつくる器官をそれぞれ答えなさい。}{精子、精巣}
\QQ{卵の核と精子の核が結びつくことを\NamedBlank{fertilization}といい、できた細胞を\NamedBlank{fertilized_egg}という。}{受精、受精卵}
\QQ{\RefBlank{fertilized_egg}から成体になるまでの過程を何というか。}{発生}

\Instruction{被子植物の有性生殖に関する次の問いに答えなさい。}
\QQ{花粉が柱頭につくことを\NamedBlank{pollination}という。}{受粉}
\QQ{\RefBlank{pollination}の後、胚珠に向かってのびる管を何というか。}{花粉管}
\QQ{被子植物の雌の生殖細胞と、それがある場所を答えなさい。}{卵細胞、胚珠の中}
\QQ{被子植物の雄の生殖細胞と、それがある場所を答えなさい。}{精細胞、花粉の中}
\QQ{被子植物では、\RefBlank{fertilization}の後、胚珠と子房はそれぞれ何になるか。}{種子、果実}

\Instruction{減数分裂と遺伝に関する次の問いに答えなさい。}
\QQ{生殖細胞ができるときに起こる特別な細胞分裂を\NamedBlank{meiosis}という。}{減数分裂}
\QQ{\RefBlank{meiosis}でできた生殖細胞の染色体数は、もとの細胞の何分の1か。}{$\dfrac{1}{2}$}
\QQ{生物のもつ形質が、親から子に伝わることを\NamedBlank{heredity}という。}{遺伝}
\QQ{形質を現すもととなるものと、その本体をそれぞれ答えなさい。}{遺伝子、DNA（デオキシリボ核酸）}
\QQ{\RefBlank{meiosis}のとき、対になった遺伝子が分かれて別々の生殖細胞に入ることを何というか。}{分離の法則}

\Instruction{遺伝の規則性に関する次の問いに答えなさい。}
\QQ{自家受粉を重ねても同じ形質が現れる個体を\NamedBlank{pure_line}という。}{純系}
\QQ{エンドウの種子の丸としわのように、どちらか一方しか現れない形質どうしを何というか。}{対立形質}
\QQ{純系の丸い種子と純系のしわのある種子を交配したとき、子に現れる形質を何というか。}{顕性形質}
\QQ{前問の交配で、子に現れない形質を何というか。}{潜性形質}
\QQ{前問の子を自家受粉させたとき、孫に現れる丸としわの数の比を答えなさい。}{丸：しわ＝3：1}
